% Part: history
% Chapter: biographies
% Section: alan-turing

\documentclass[../../../include/open-logic-section]{subfiles}

\begin{document}

\olfileid{his}{bio}{tur}

\olsection{অ্যালান টুরিং}

অ্যালান টুরিংয়ের জন্ম ১৯১২ সালের ২৩ জুন
লন্ডনের মাইডা ভেলে। তাঁকে তাত্ত্বিক কম্পিউটার
বিজ্ঞানের জনক মনে করা হয়। ছোটবেলা থেকেই
প্রাকৃতিক বিজ্ঞান ও গণিতে টুরিংয়ের আগ্রহ ছিল।
কিন্তু বালকবয়সে তিনি যেসব বিদ্যালয়ে
পড়তেন, সেখানে তাঁর আগ্রহের বিষয়গুলির
যথেষ্ট গুরুত্ব ছিল না; সাহিত্য ও ধ্রুপদি
বিদ্যার উপরেই বেশি জোর দেওয়া হতো।
ফলে বিদ্যালয়ে তাঁর ফল ভালো হয়নি এবং
অনেক শিক্ষক তাঁকে তিরস্কার করতেন।

\olphoto{turing-alan}{অ্যালান টুরিং}

টুরিং কেমব্রিজের কিংস কলেজে স্নাতক
পর্যায়ে গণিত পড়েন। গণনসাধ্য অপেক্ষকের
ধারণাটি নির্ভুলভাবে সংজ্ঞায়িত করা এবং
সিদ্ধান্ত-সমস্যার অসিদ্ধান্তযোগ্যতা প্রমাণ
করার চেষ্টায় ১৯৩৬ সালে তিনি বর্তমানে
টুরিং যন্ত্র নামে পরিচিত ধারণাটি নির্মাণ
করেন। অ্যালোনজো চার্চ নিজের ল্যাম্বডা
ক্যালকুলাসের সাহায্যে ফলটি আগেই
প্রমাণ করে ফেলেছিলেন। তা সত্ত্বেও
চার্চের ফলের উল্লেখসহ টুরিংয়ের
প্রবন্ধটি প্রকাশিত হয়। চার্চ তাঁকে
প্রিন্সটনে আমন্ত্রণ জানান; টুরিং
সেখানে ১৯৩৬--১৯৩৮ সাল কাটান
এবং চার্চের অধীনে ডক্টরেট লাভ করেন।

যুক্তিবিদ্যায় আগ্রহী হলেও প্রাকৃতিক
বিজ্ঞানের প্রতি টুরিংয়ের পুরোনো
আগ্রহ বজায় ছিল। দ্বিতীয় বিশ্বযুদ্ধের
সময় ব্লেচলি পার্কে ব্রিটিশ সংকেতবিশ্লেষণ
বিভাগে কাজ করার সময় তাঁর ব্যবহারিক
দক্ষতা কাজে লাগে। জার্মান নৌবাহিনীর
যোগাযোগে ব্যবহৃত এনিগমা সংকেত
ভাঙার কাজে টুরিং ছিলেন কেন্দ্রীয়
ব্যক্তিত্ব। পরিসংখ্যান ও সংকেতলিপিবিদ্যায়
তাঁর দক্ষতা এবং তড়িৎযান্ত্রিক যন্ত্রের
ব্যবহারের ফলে দলটি “bombe” নামে
একটি সংকেত-পাঠোদ্ধারকারী যন্ত্র
নির্মাণ করে সংকেতটি ভাঙতে পারে।
তাঁর ধারণাগুলি পৃথিবীর প্রথম
প্রোগ্রামযোগ্য বৈদ্যুতিন কম্পিউটার
Colossus নির্মাণেও সাহায্য করে;
ব্লেচলি পার্কে সেটি জার্মান লোরেঞ্জ
সংকেত ভাঙতে ব্যবহৃত হয়।
% BN-SRC-853 TARGET-CORRECTION-BEGIN
\par\noindent\textbf{পাঠ-সংশোধন।} Bombe ছিল তড়িৎযান্ত্রিক যন্ত্র; মূলপাঠের
বৈদ্যুতিন বিশেষণটি সংশোধিত। দেখো
\href{https://www.gchq.gov.uk/person/peter-twinn}{জিসিএইচকিউ-র টুইন-সংক্রান্ত নথি}।
% BN-SRC-853 TARGET-CORRECTION-END

টুরিং ছিলেন সমকামী। তা সত্ত্বেও
১৯৪১ সালে ব্লেচলি পার্কের সহকর্মী
জোয়ান ক্লার্ককে তিনি বিয়ের প্রস্তাব
দেন; প্রস্তাবের পরেই তাঁকে নিজের
সমকামিতার কথা জানান। পরে বাগদান ভেঙে দেন। জীবনে তাঁর
কয়েকজন প্রেমিক ছিলেন, যদিও সে
সময়ে যুক্তরাজ্যে সমকামী যৌনসম্পর্ক
ফৌজদারি অপরাধ ছিল। ১৯৫২ সালে
তাঁর তৎকালীন প্রেমিকের এক বন্ধু
টুরিংয়ের বাড়িতে চুরি করে। পুলিশে
অভিযোগ জানাতে গিয়ে টুরিং
সমকামী সম্পর্ক থাকার কথা স্বীকার
করেন; তিনি ভেবেছিলেন সরকার
শিগগিরই সমকামী যৌনসম্পর্ককে
বৈধ করতে চলেছে। তা সত্য ছিল না।
তাঁর বিরুদ্ধে গুরুতর অশালীনতার
অভিযোগ আনা হয়। কারাগারে
যাওয়ার বদলে টুরিং যৌনাকাঙ্ক্ষা
কমানোর হরমোন চিকিৎসা বেছে
নেন। ১৯৫৪ সালের ৮ জুন
সায়ানাইডের অতিমাত্রার কারণে
তাঁকে মৃত অবস্থায় পাওয়া যায়—খুব
সম্ভবত এটি আত্মহত্যা ছিল।
২০১৩ সালে রানি দ্বিতীয় এলিজাবেথ
তাঁকে রাজকীয় ক্ষমা দেন।
% BN-SRC-854 TARGET-CORRECTION-BEGIN
\par\noindent\textbf{পাঠ-সংশোধন।} বাগদান ১৯৪১ সালের, ১৯৪২ নয়; প্রস্তাবের পরেই
নিজের সমকামিতার কথা জানানো ও পরে বাগদান ভাঙার ক্রমও সংশোধিত।
দেখো \href{https://www.gchq.gov.uk/information/joan-clarke}{জিসিএইচকিউ-র ক্লার্ক-সংক্রান্ত নথি} ও
\href{https://www.turing.org.uk/scrapbook/ww2.html}{হজেসের বিবরণ}।
% BN-SRC-854 TARGET-CORRECTION-END

\begin{reading}
অ্যালান টুরিংয়ের পূর্ণাঙ্গ জীবনী
জানতে \citet{Hodges2014} দেখো।
টুরিংয়ের জীবন ও কাজ থেকে
অনুপ্রাণিত \emph{Breaking the Code}
নাটকটি ১৯৯৬ সালে টেলিভিশনের
জন্য নির্মিত হয়; ডেরেক জ্যাকবি
টুরিংয়ের ভূমিকায় অভিনয় করেন।
বেনেডিক্ট কাম্বারব্যাচ ও কিয়েরা
নাইটলি অভিনীত, অ্যাকাডেমি
পুরস্কারের জন্য মনোনীত
\emph{The Imitation Game}
চলচ্চিত্রটিও টুরিংয়ের জীবন ও
ব্লেচলি পার্কের সময়কে আলগাভাবে
অবলম্বন করেছে \citep{Imitation2014}।

টুরিংয়ের জীবন ও কাজ নিয়ে
\citet{Radiolab2012}-এ কয়েকটি
পডকাস্ট আছে। বিবিসি হরাইজনের
\emph{The Strange Life and Death of
Dr.~Turing} প্রামাণ্যচিত্রটি অনলাইনে
দেখা যায় \citep{Sykes1992}।
\citep{Theelen2012}-এ সচল একটি
LEGO টুরিং যন্ত্রের সংক্ষিপ্ত ভিডিও
আছে; ২০১২ সালে টুরিংয়ের জন্মশতবর্ষ
উপলক্ষে সেটি তৈরি হয়েছিল।

টুরিং যন্ত্র ও সিদ্ধান্ত-সমস্যা নিয়ে
টুরিংয়ের মূল প্রবন্ধটি হলো
\citet{Turing1937}।
\end{reading}

\end{document}
